\documentclass[11pt,a4paper]{article}
\usepackage[T1]{fontenc}
\usepackage{lmodern}
\usepackage[margin=25mm]{geometry}
\usepackage{amsmath,amssymb}
\usepackage{booktabs,tabularx}
\usepackage[hidelinks]{hyperref}
\usepackage{microtype}
\setlength{\parindent}{0pt}
\setlength{\parskip}{0.5em}
\clubpenalty=10000
\widowpenalty=10000
\emergencystretch=2em
\newcommand{\E}{\mathbb{E}}
\newcommand{\code}[1]{\texttt{#1}}

\begin{document}
{\LARGE\bfseries PLDSC: univariate and bivariate LD Score regression\par}
\medskip
{\small Methods and implementation notes, 29 October 2026.\\
Reviewed against the PLDSC 0.1.0.dev1 source tree on that date.\par}

\medskip
PLDSC estimates SNP heritability from one GWAS and genetic correlation from
two GWAS using LD Score regression (LDSC). It also contains research tools
for assessing LD-score mismatch in multi-ancestry meta-analyses. This report
explains the ordinary LDSC fits, their inputs and computational design. The
research results remain in \code{docs/MULTI\_ANCESTRY.md}; the runnable
examples and full API guide are in \code{docs/LDSC.md}.

\section{Inputs and targets}

For variant $j$ in trait $t$, let $z_{tj}=b_{tj}/s_{tj}$ be the signed GWAS
effect divided by its standard error, and let $N_{tj}$ be the appropriate
effective sample size. For a reference LD correlation $r_{jk}$, define the
LD score over a reference variant set $\mathcal M$ of size $M$ as
\begin{equation}
 \ell_j=\sum_{k\in\mathcal M} r_{jk}^{2}.
 \label{eq:ldscore}
\end{equation}
The reference must match the populations and genome build of the GWAS well
enough for the intended analysis, and the score for each regression variant
must be associated with the correct reference row. $M$ counts the variants
summed into the score; it is not, in general, the number of GWAS rows left
after QC. PLDSC fits one common LD-score vector for the two traits, so a
single reference must be appropriate for both.

The targets are the SNP heritabilities $h_1^2,h_2^2$ and the additive
genetic covariance $g_{12}$ for the modeled variant set and scale. Their
genetic correlation is
\begin{equation}
 r_g=\frac{g_{12}}{\sqrt{h_1^2 h_2^2}}.
 \label{eq:rg}
\end{equation}
This is a correlation of genetic components, not a phenotypic correlation or
a causal effect. A binary trait is not automatically transformed to the
liability scale. Reference mismatch, effect-size architecture, selection and
imperfect LD estimates can alter the interpretation of the fitted slope.

\section{The regressions}

Under the standard LDSC approximation, with standardized effects having
roughly constant variance and covariance across variants, the conditional
moments follow Equations~\eqref{eq:univ} and~\eqref{eq:cross}
\cite{ldsc,rgldsc}:
\begin{align}
 \E[z_{tj}^{2}\mid\ell_j] &\approx a_t + \frac{N_{tj}\ell_j}{M} h_t^2,
 \label{eq:univ}\\
 \E[z_{1j}z_{2j}\mid\ell_j] &\approx a_{12} +
   \frac{\sqrt{N_{1j}N_{2j}}\,\ell_j}{M} g_{12}.
 \label{eq:cross}
\end{align}
The slope of Equation~\eqref{eq:univ} estimates $h_t^2$; the slope of
Equation~\eqref{eq:cross} estimates $g_{12}$. PLDSC fits both marginal
regressions and the cross-trait regression on the \emph{same common rows}
for a bivariate analysis, and then uses Equation~\eqref{eq:rg}. A separate
one-trait fit can have a different $h^2$ because it can retain different
variants. The two GWAS effects must count the same allele before their signed
$z$-scores are multiplied.

The intercepts are free by default. The marginal intercept can reflect
confounding, genomic-control effects or model/reference mismatch; the
cross-trait intercept can also reflect sample overlap and correlated
confounding. These sources are not identified separately by the intercept.
Fixing the intercept is an assumption, not a diagnostic: in the bivariate
fit \code{constrain\_intercept} fixes only $a_{12}$, leaving $a_1,a_2$ free.
A misspecified fixed intercept can bias the slopes.

\subsection{Weights and uncertainty}

PLDSC uses one-step iteratively reweighted least squares. Let
$x_{tj}=N_{tj}\ell_j/M$ and $L_j=\max(\ell_j,1)$. A marginal fit starts with
$\mu_{tj}=1$ and, after each fit, replaces it by
$\max(a_t+h_t^2x_{tj},1)$. Its weights are
\begin{equation}
 w_{tj}=\frac{1}{2\mu_{tj}^{2}L_j}.
 \label{eq:marginalweight}
\end{equation}
For the cross-trait fit, the fitted marginal means give the approximate
bivariate-normal variance of $z_{1j}z_{2j}$. With fitted cross mean $q_j$,
PLDSC uses
\begin{equation}
 w_{12,j}=\frac{1}{L_j\max(\mu_{1j}\mu_{2j}+q_j^2,10^{-6})}.
 \label{eq:crossweight}
\end{equation}
The first cross fit uses $q_j=0$ when its intercept is free, or the fixed
intercept otherwise; subsequent fits update $q_j$. By default there are two
weight updates after the initial fit. This is a one-step estimator: PLDSC
does not reproduce every option, including the original software's two-step
intercept estimator, of the standalone LDSC implementation.

Standard errors use a delete-a-block jackknife over contiguous variants in
genomic order. If $\widehat\theta_{(-b)}$ is the estimate after deleting
block $b$, with mean $\overline\theta_{(-\cdot)}$ across $B$ deletions,
\begin{equation}
 \widehat{\operatorname{se}}(\widehat\theta)=
 \left[\frac{B-1}{B}\sum_{b=1}^{B}
 (\widehat\theta_{(-b)}-\overline\theta_{(-\cdot)})^2\right]^{1/2}.
 \label{eq:jackknife}
\end{equation}
Each deletion refits the weights; a bivariate deletion refits both marginals
and the cross-trait regression before calculating $r_g$. The defaults are
100 blocks for one trait and 200 for two, capped by the number of regression
variants. For $r_g$, the point estimate is undefined if either marginal
estimate is non-positive; its standard error is undefined if a required
replicate is invalid or fewer than two blocks are available. Estimates are
not clipped to their theoretical ranges.

\section{From files to estimates}

Figure~\ref{fig:workflow} summarizes the file workflow. Allelign reads each
GWAS, checks its columns and sample sizes, filters rows, and harmonizes
alleles to a sorted \code{.bim} or \code{.pvar} reference. Defaults remove
non-finite effects, non-positive standard errors or sample sizes, duplicated
identifiers, and variants with $N$ below 70\% of the median. Where available,
MAF below 0.01 and INFO below 0.7 are filtered. Ambiguous A/T and C/G sites
are dropped by default, and effect signs are flipped where alignment requires
it. QC and alignment counts are returned for each trait.

\begin{figure}[htbp]
\centering\small
\setlength{\fboxsep}{4pt}
\fbox{\parbox[c][2.0cm][c]{0.23\linewidth}{\centering GWAS summaries\\Reference variants\\LD scores and $M$}}
\enspace$\longrightarrow$\enspace
\fbox{\parbox[c][2.0cm][c]{0.23\linewidth}{\centering Allelign QC\\Allele alignment\\Chi-square cap}}
\enspace$\longrightarrow$\enspace
\fbox{\parbox[c][2.0cm][c]{0.23\linewidth}{\centering LDSC fits\\Block jackknife\\Estimates and logs}}
\caption{The file-based workflow for one or two traits. Bivariate fits use
common, aligned variants passing the cap for both traits.}
\label{fig:workflow}
\end{figure}

The wrappers then apply the per-trait cap $z_{tj}^{2}\leq
\max(80,0.001N_{tj})$; a bivariate row must pass both caps. They do not
automatically remove long-range LD regions or check genome builds. Such
selection, liftover and reference suitability are the analyst's responsibility.
A numerical \code{n\_eff} passed as a read option replaces every per-variant
sample size. For a case-control study, a common choice is
\begin{equation}
 N_{\rm eff}=\frac{4}{1/N_{\rm case}+1/N_{\rm control}}.
 \label{eq:neff}
\end{equation}
For compatible meta-analyses of non-overlapping cohorts, the sum of
cohort-specific effective sizes can differ from applying Equation~\eqref{eq:neff}
to pooled counts; choose the scale appropriate to the GWAS.

Table~\ref{tab:api} separates the file-based and preprocessed-array paths.
The latter do not perform QC, harmonization, filtering or sorting. The
\code{ldsc\_rg} array function takes LDpred3-standardized effects; raw
per-allele GWAS effects belong in the file wrapper or must first be converted
to signed $z$-scores for \code{ldsc\_rg\_z}. For example, with file paths
\code{gwas1}, \code{gwas2}, and \code{ref}, and reference-aligned
\code{scores}, the file-based calls are:
{\small
\begin{verbatim}
h2 = ldsc_h2_sumstats(gwas1, ref, scores, m_snps=M)
rg = ldsc_rg_sumstats(gwas1, gwas2, ref, scores,
                      m_snps=M, ncores=4)
print(h2.estimate.h2, rg.estimate.rg, rg.estimate.rg_se)
\end{verbatim}
}
Both functions are imported from \code{pldsc}. Their results also contain QC
and alignment logs; the complete, runnable example is in the guide.

\begin{table}[htbp]
\centering\small
\caption{Main entry points and their inputs. The file wrappers return QC
and alignment logs alongside the estimate.}
\label{tab:api}
\begin{tabularx}{\linewidth}{@{}lX@{}}
\toprule
Function & Input and result\\ \midrule
\code{ldsc\_h2\_sumstats} & One GWAS and reference: SNP heritability and SE.\\
\code{ldsc\_rg\_sumstats} & Two GWAS and reference: genetic covariance,
correlation and correlation SE.\\
\code{ldsc\_h2} & Aligned arrays of $z^2$, LD scores and sample sizes.\\
\code{ldsc\_rg\_z} & Aligned arrays of signed $z$-scores, LD scores and sample sizes.\\
\bottomrule
\end{tabularx}
\end{table}

\section{Scores, computation and evidence}

\code{ld\_scores} accepts a dense correlation matrix or disjoint indexed
blocks, including int8 and low-rank representations. Low-rank contractions
avoid materializing the full dense matrix. Block scores omit between-block
LD, and the optional finite-reference correction can yield non-positive
scores. The regressions require positive scores; silently clipping them
changes the analysis. Score construction and the choice of $M$ must match
the intended reference, window and variant set.

The regression kernels were adapted from LDpred3 and bipred and are designed
for very fast repeated fits. With optional Numba compilation, delete-block fits traverse retained ranges without making
a full copy for each deletion. Independent bivariate deletions can run in
parallel; poorly conditioned fits fall back to the reference numerical path.
Without Numba, a NumPy implementation remains available. Compilation affects
the first call, while elapsed time and memory depend on the workload and
hardware. The repository includes a paired warm timing driver; no quantitative
speedup is asserted here without a corresponding measured run.

\pagebreak
The tests include an exact univariate regression line, dense and low-rank
score checks, near-singular fits, allele flips and file-workflow checks, and
comparisons with LDpred3 and bipred when those siblings are installed. These
are implementation and consistency checks, not evidence of real-data
calibration. The synthetic example exercises the API but does not model LD
among GWAS statistics.

\section{Relation to the population-structure research}

The ordinary regressions above use one LD-score vector. The multi-ancestry
research instead derives a meta-analysis score containing within- and
cross-ancestry LD and examines the consequences of using a single-ancestry
reference. Its ancestry-pair effect correlation is not the cross-trait $r_g$
in Equation~\eqref{eq:rg}. In the current panel experiments, reported
$\kappa$ values are panel plug-in slope ratios under stated assumptions;
finite-reference noise can affect both scores, so these are not calibrated
estimates of absolute heritability bias. The definitions, assumptions,
experiments and results are recorded in \code{docs/MULTI\_ANCESTRY.md}.

\begin{thebibliography}{9}
\bibitem{ldsc} Bulik-Sullivan BK, Loh PR, Finucane HK, et al.
LD Score regression distinguishes confounding from polygenicity in
 genome-wide association studies. \emph{Nature Genetics} 47, 291--295 (2015).
\href{https://doi.org/10.1038/ng.3211}{doi:10.1038/ng.3211}.
\bibitem{rgldsc} Bulik-Sullivan B, Finucane HK, Anttila V, et al.
An atlas of genetic correlations across human diseases and traits.
\emph{Nature Genetics} 47, 1236--1241 (2015).
\href{https://doi.org/10.1038/ng.3406}{doi:10.1038/ng.3406}.
\end{thebibliography}
\end{document}
